% Working draft for a Journal of Open Source Software submission.
% The manuscript describes Biogeme 3.3.5. Replace the version and archive
% information below if the submission is made against a later release.
\documentclass[11pt]{article}

\usepackage[T1]{fontenc}
\usepackage[utf8]{inputenc}
\usepackage{lmodern}
\usepackage[margin=1in]{geometry}
\usepackage{microtype}
\usepackage{booktabs}
\usepackage{array}
\usepackage{tabularx}
\usepackage{url}
\usepackage[hidelinks]{hyperref}
\usepackage[numbers,sort&compress]{natbib}

\newcolumntype{Y}{>{\raggedright\arraybackslash}X}

\title{Biogeme: An open-source Python framework for research and education in discrete-choice modeling}

\author{Michel Bierlaire\\
  Transport and Mobility Laboratory, ENAC,\\
  \textsc{École Polytechnique Fédérale de Lausanne}, Switzerland}

\date{}

\begin{document}

\maketitle

\begin{abstract}
Biogeme is an open-source Python package for the estimation and application of
parametric models, with a particular emphasis on discrete-choice analysis. It
is designed both as a research framework for developing and estimating advanced
models and as an educational platform through which students can learn model
specification, simulation, estimation, and interpretation. Its symbolic
expression system allows users to write transparent likelihoods and reuse the
same specification for estimation, simulation, derivatives, validation, and
reporting. The package supports multinomial, nested, cross-nested, mixed,
latent-class, ordered, panel, sampled-alternative, MDCEV, and hybrid choice
models, together with maximum-likelihood and Bayesian workflows. NumPy and JAX
backends, Monte Carlo integration, numerical-safety controls, memory-aware
derivatives, and recoverable results make the framework suitable for demanding
empirical applications while retaining a readable interface for teaching.
\end{abstract}

\section{Summary}

Biogeme is an open-source Python package for the maximum-likelihood estimation
of general parametric models, with a long-standing focus on discrete-choice
models. The project grew from the original BIOGEME software described by
\citet{Bierlaire2003,BierlaireFetiarison2009} and has evolved into a framework
for specifying, estimating, simulating, and documenting a broad range of
behavioral models. The current release considered in this paper is Biogeme
3.3.5.

The package has two closely related goals. For researchers, it provides a
general expression language in which standard and non-standard likelihoods can
be written directly, combined with simulation, automatic differentiation,
Bayesian inference, and detailed post-estimation analysis. For educators, it
provides a consistent syntax and a large collection of runnable examples that
make the connection between equations, code, estimation output, and behavioral
interpretation explicit. The same interface can therefore support an
introductory multinomial logit exercise and a research application involving
latent variables, panel observations, or a cross-nested structure.

Biogeme is distributed under the MIT License. The source code is available at
\url{https://github.com/michelbierlaire/biogeme}, the package is distributed on
PyPI, and the documentation is hosted at
\url{https://biogeme.epfl.ch/sphinx/index.html}.

\section{Statement of need}

Discrete-choice models are used to study decisions among competing
alternatives in transportation, marketing, environmental economics, health,
and many other fields. Although the multinomial logit model is an important
starting point, applied research quickly encounters richer structures:
correlated alternatives, random coefficients, repeated choices, latent
preferences, endogenous sampling, multiple discrete-continuous decisions, and
user-defined likelihood contributions. Estimating such models requires more
than a collection of isolated estimators. It requires a representation that is
flexible enough to express the model, reliable numerical methods, and tools for
examining the implications of the fitted model.

Many software packages provide excellent implementations of particular model
families. A researcher who needs a model outside those families, however, may
have to derive and implement likelihoods, gradients, simulation, and reporting
separately. This separation makes experimentation difficult and can make it
hard for readers or students to see how the implemented model corresponds to
the mathematical specification. Biogeme addresses this need by treating the
model expression as a first-class object that is shared by the computational
workflows.

The educational need is related but distinct. Students should be able to start
from a small, readable example and then extend it without learning a new
software interface for each model. A teaching package should expose the
structure of the likelihood rather than hide it behind a fixed formula, while
still providing robust defaults and informative output. Biogeme's examples
cover a progression from basic logit and binary models through simulation,
mixtures, GEV models, panel data, Bayesian estimation, and hybrid choice
models. This makes the package useful in courses as well as in research
projects.

\section{State of the field}

Biogeme occupies a space shared by general statistical environments, packages
specialized in random-utility models, and commercial econometric software. The
following comparison is intentionally high-level: it identifies design
emphases rather than attempting to rank packages or reproduce every feature
of their current releases.

\begin{table}[htbp]
  \centering
  \caption{Representative software related to Biogeme.}
  \label{tab:state-field}
  \begin{tabularx}{\textwidth}{@{}lYY@{}}
    \toprule
    Software & Primary orientation & Relation to Biogeme in this paper \\
    \midrule
    Biogeme & Open-source Python framework with symbolic expressions,
      simulation, advanced discrete-choice structures, and Bayesian workflows. &
      Research and education through one extensible specification language. \\
    \texttt{mlogit} & R package for maximum-likelihood random-utility models,
      including several standard extensions \citep{Croissant2020}. &
      Strong specialized alternative for common random-utility models. \\
    \texttt{gmnl} & R package focused on generalized and mixed multinomial
      logit models \citep{SarriasDaziano2017}. &
      Focused implementation for important classes of taste heterogeneity. \\
    Apollo & R framework for classical and Bayesian choice-model estimation,
      including user-defined model components \citep{HessPalma2019}. &
      Broad choice-model environment with a different language and workflow. \\
    \texttt{statsmodels} & General Python statistical-modeling library whose
      discrete module includes multinomial logit and related models
      \citep{SeaboldPerktold2010}. &
      General statistical ecosystem rather than a dedicated choice framework. \\
    NLOGIT & Proprietary econometric software for multinomial, mixed, nested,
      latent-class, and related discrete-choice models \citep{NLOGIT2025}. &
      Commercial benchmark with a long history in applied econometrics. \\
    ALOGIT & Proprietary professional software for estimating and analyzing
      generalized logit choice models \citep{ALOGIT2025}. &
      Commercial software emphasizing large logit-family applications. \\
    \bottomrule
  \end{tabularx}
\end{table}

Biogeme's distinctive design choice is to combine a general symbolic
representation with dedicated implementations of important choice structures.
The result is not intended to replace every specialized package. Rather, it
offers a transparent route from a mathematical expression to estimation,
simulation, derivatives, and reports, including when the expression does not
belong to a small catalogue of predefined models.

\section{Software design}

\subsection{Symbolic expressions as the common model representation}

Biogeme represents variables, parameters, constants, transformations, and
probability expressions as nodes in an expression tree. Users can combine these
objects with ordinary arithmetic notation to construct utilities, probabilities,
measurement equations, structural equations, and complete likelihoods. The
tree can then be evaluated for data and parameter values, differentiated, or
translated to a numerical backend. This avoids maintaining separate versions
of the model for estimation and simulation and makes the source code close to
the mathematical specification.

The expression representation also supports generic likelihoods. A user can
define a contribution that is not one of the built-in choice probabilities and
still use the standard estimation, derivative, reporting, and validation
machinery. At the same time, dedicated expressions are available for common
structures where exploiting the model's algebra improves clarity or
performance.

\subsection{Choice-model structures}

The standard example sequence includes multinomial and binary logit, nested and
cross-nested logit, mixed and latent-class models, ordered logit and probit,
panel likelihoods, sampled alternatives, and models with individual-level
parameters. Monte Carlo integration can use pseudo-random or quasi-random draw
schemes, including Halton and modified Latin hypercube sequences, symmetric
draws, and antithetic constructions. The simulation framework is shared by
estimation and post-estimation calculations.

Biogeme also supports multiple-discrete continuous extreme-value (MDCEV)
models, including applications with no outside good \citep{BierlaireWang2024},
and estimation with samples of alternatives \citep{BierlairePaschalidis2023}.
Assisted specification tools help construct and compare families of models,
while parameter overrides allow a user to control automatically generated
parameters without rewriting the surrounding specification.

\subsection{Hybrid choice models}

Hybrid choice models combine a discrete-choice component with latent variables,
structural equations, and measurement equations. Biogeme provides objects for
these components, normalization rules, ordinal and continuous measurements,
and both sequential and simultaneous estimation. The framework described in
\citet{Bierlaire2016Latent,BierlaireBenAkivaWalker2026} is designed to keep the
behavioral interpretation visible in the model code while allowing the
measurement system and the choice system to be estimated jointly.

\subsection{Numerical backends and derivatives}

The conventional evaluation path uses NumPy and is intended to be easy to
inspect and debug. JAX can compile expression evaluations and provide automatic
differentiation for gradients and higher-order derivatives. Biogeme exposes
numerical-safety controls because the most appropriate trade-off depends on the
model and data. In fast mode, expressions use direct arithmetic whenever this
is appropriate. In numerically safe mode, guarded and log-domain operations
are used to reduce avoidable overflow, underflow, and invalid intermediate
values.

Several implementation choices target models whose dimensions make derivatives
expensive. Sparse cross-nested logit expressions avoid work for alternatives
with literal zero allocations. Hessians can be evaluated in chunks when a full
automatic-differentiation workload would exceed a practical memory budget.
These options preserve the model expression while making the computational
strategy explicit to the user.

\subsection{Estimation, diagnostics, and results}

Biogeme supports maximum-likelihood estimation with several optimization
strategies and derivative configurations. It can compute gradients, Hessians,
and BHHH matrices, perform bootstrap and validation calculations, and produce
elasticities, willingness-to-pay measures, and simulated quantities. Bayesian
workflows are available through the Python probabilistic-programming
ecosystem.

Estimation results are saved in YAML files that can be recovered after an
interruption, while HTML, LaTeX, pandas, and NetCDF outputs support inspection
and downstream analysis. The separation between recoverable machine-readable
results and human-readable reports is useful both in long research runs and in
teaching, where students can inspect each stage of an estimation exercise.

\subsection{Research and educational design}

The example collection is organized as a progression. Early examples introduce
data, utilities, availability, weights, and standard logit estimation. Later
examples add simulation, mixtures, GEV structures, panel observations,
multiple models, ordinal measurements, Bayesian estimation, and latent
variables. Each example is executable and is intended to show a complete
workflow rather than only an isolated API call.

This progression is also a software-design constraint: model definitions should
remain readable when new features are added. A student can modify a utility or
probability expression and rerun the same workflow, while a researcher can use
the same expression tree as the starting point for a more specialized model.

\section{Research and educational impact}

\subsection{Research impact}

Biogeme has been developed and publicly distributed over several decades. Its
technical reports document successive extensions for Monte Carlo integration
\citep{Bierlaire2015}, latent-variable models \citep{Bierlaire2016Latent},
assisted specification \citep{Ortelli2020}, sampled alternatives
\citep{BierlairePaschalidis2023}, and MDCEV models \citep{BierlaireWang2024}.
The package has also been used in empirical studies of transportation,
connected vehicles, information effects, electric mobility, and other
applications. Representative examples include \citet{Chorus2007},
\citet{Derikx2016}, \citet{Kajanova2021}, \citet{Gerzinic2022}, and
\citet{Nigro2024}.

These applications illustrate the role of Biogeme as a research instrument:
the software is used not only to estimate standard models, but also to explore
heterogeneity, latent preferences, policy scenarios, and new likelihood
specifications. The public repository, issue tracker, release history, and
documentation make the computational environment inspectable by other
researchers.

\subsection{Educational impact}

Biogeme is also intended for teaching. Its examples can be used to introduce
the logic of random-utility models, maximum likelihood, simulated likelihood,
identification, parameter constraints, and post-estimation interpretation.
The progression from small examples to hybrid choice models supports courses
with different levels of mathematical and computational depth. Because the
model expression is written explicitly, instructors can connect a lecture's
equations to executable code without introducing a separate opaque model
language.

The open-source license and public documentation lower the barrier to classroom
use. Students can inspect the implementation, reproduce published examples,
and extend them for projects. The same scripts can then be retained as a
starting point for research, making the transition from education to applied
work direct.

\section{Reproducibility and quality assurance}

The repository contains a large automated test suite covering expressions,
databases, estimation, simulation, reporting, Bayesian workflows, and example
infrastructure. Continuous-integration workflows exercise supported Python
versions on Linux, macOS, and Windows, and documentation builds include link
and example checks. Releases are built from versioned source, tested, and
published through a reproducible packaging workflow.

The examples use public or repository-provided data and record their outputs in
machine-readable and human-readable formats. Recoverable optimization results
allow long estimations to resume without discarding completed work. For
simulation-based models, the configuration of draws and numerical options is
part of the model workflow, so users can report the computational choices that
affect simulation noise and run time.

\section{Availability and sustainability}

Biogeme is available from PyPI and from its public GitHub repository. The
repository includes installation instructions, documentation, a comprehensive
example gallery, tests, and release procedures. Questions and user support are
handled through the project's public communication channels. The software is
distributed under the MIT License, allowing use, modification, and
redistribution with attribution.

The project is maintained by the author at the Transport and Mobility
Laboratory of EPFL. Development and maintenance are funded by the laboratory;
there was no external funder or sponsor influence on the design, implementation,
testing, or release decisions described here. A permanent archival record and
DOI for the submitted release will be added to this section before submission.

\section{Conclusions}

Biogeme combines a transparent symbolic specification language with numerical
and reporting infrastructure for discrete-choice research. Its coverage of
standard and advanced model structures, simulation, latent variables, Bayesian
workflows, and post-estimation analysis makes it suitable for methodological
development and demanding empirical applications. Its readable syntax, public
examples, and progressive documentation also make it suitable for education.

The central design principle is continuity: the expression introduced by a
researcher or instructor is the object that is evaluated, differentiated,
estimated, simulated, and reported. This reduces duplication between a model's
mathematical description and its implementation, while leaving users in control
of the computational choices that matter for speed, numerical stability, and
memory use.

\section*{AI usage disclosure}

ChatGPT 5.6 was used on September 2, 2026, solely to assist with the
organization, drafting, and editing of this manuscript. It was not used for
software development, implementation, testing, or release decisions. Michel
Bierlaire reviewed and validated the manuscript, made all substantive technical
decisions, and accepts full responsibility for its content.

\section*{Acknowledgements}

The author acknowledges the support of the Transport and Mobility Laboratory at
the École Polytechnique Fédérale de Lausanne.

\bibliographystyle{plainnat}
\bibliography{biogeme_joss}

\end{document}
